% Part: first-order-logic
% Chapter: model-theory
% Section: reducts-and-expansions

\documentclass[../../../include/open-logic-section]{subfiles}

\begin{document}

\olfileid{mod}{bas}{red}
\section{فروکاست‌ها و بسط‌ها}

اغلب مقایسهٔ زبان‌هایی که نمادهای مشترک دارند، و نیز مقایسهٔ
!!{structure}s این زبان‌ها، سودمند یا ضروری است. رایج‌ترین حالت زمانی
است که همهٔ نمادهای !!a{language}~$\Lang{L}$ بخشی از !!a{language}
چون~$\Lang{L'}$ نیز باشند، یعنی $\Lang{L} \subseteq \Lang{L'}$.
آنگاه همواره می‌توان یک !!{structure}-$\Lang{L}$ چون~$\Struct{M}$ را
با افزودن تفسیرهایی برای نمادهای افزوده، در حالی که تفسیر نمادهای مشترک
یکسان می‌ماند، به یک !!{structure}-$\Lang{L'}$ بسط داد. از سوی دیگر،
می‌توانیم از ساختارِ~$\Lang{L'}$ چون~$\Struct{M'}$ صرفاً با
«فراموش‌کردن» تفسیر نمادهایی که در~$\Lang{L}$ رخ نمی‌دهند، یک
ساختاری برای~$\Lang{L}$ به دست آوریم.

\begin{defn}
\ollabel{defn:reduct}
فرض کنید $\Lang L \subseteq \Lang L'$، $\Struct M$ یک
!!{structure}-$\Lang L$ و $\Struct M'$ یک
!!{structure}-$\Lang L'$ باشد. $\Struct M$
\emph{فروکاستِ}~$\Struct M'$ به~$\Lang L$، و $\Struct M'$
\emph{بسطِ}~$\Struct M$ به~$\Lang L'$ است هرگاه و تنها هرگاه
\begin{enumerate}
\item $\Domain{M} = \Domain{M'}$
\item برای هر !!{constant}~$c \in \Lang L$، $\Assign{c}{M} =
  \Assign{c}{M'}$.
\item برای هر !!{function}~$f \in \Lang L$، $\Assign{f}{M} =
  \Assign{f}{M'}$.
\item برای هر !!{predicate}~$P \in \Lang L$، $\Assign{P}{M} =
  \Assign{P}{M'}$.
\end{enumerate}
\end{defn}

\begin{prop}
\ollabel{prop:reduct}
اگر !!{structure}-$\Lang{L}$ چون~$\Struct{M}$ فروکاستِ یک
!!{structure}-$\Lang{L'}$ چون~$\Struct{M'}$ باشد، آنگاه برای همهٔ
!!{sentence}s-$\Lang{L}$ چون~$!A$،
\[
\Sat{M}{!A} \text{ هرگاه و تنها هرگاه } \Sat{M'}{!A}.
\]
\end{prop}

\begin{proof}
  تمرین.
\end{proof}

\begin{prob}
\olref[mod][bas][red]{prop:reduct} را اثبات کنید.
\end{prob}

\begin{defn}
فرض کنید یک ساختارِ~$\Lang{L}$ چون~$\Struct{M}$ داریم، و
$\Lang{L'} = \Lang{L} \cup \{P\}$ بسطِ~$\Lang{L}$ باشد که با افزودن
یک !!{predicate}~$n$موضعی چون~$P$ به دست آمده است، و
$R \subseteq \Domain{M}^n$ رابطه‌ای $n$موضعی باشد. در این صورت، برای
بسط~$\Struct{M'}$ از~$\Struct{M}$ که $\Assign{P}{M'} = R$، می‌نویسیم
$\Expan{M}{R}$.
\end{defn}

\end{document}
